\documentclass{article}
\usepackage{xltxtra}
\usepackage{fontspec}
\usepackage{xeCJK}

% Inscription macros matching original book format
\newcommand{\inscriptionref}[1]{\textit{(#1)}}
\newcommand{\inscriptionlabel}[1]{\textbf{(#1:)}}
\newcommand{\inscriptionside}[1]{\emph{#1}}
\newcommand{\inscriptionsection}[1]{%
  \par\vspace{0.5\baselineskip}
  \noindent\textbf{\MakeUppercase{#1}}%
  \par\vspace{0.3\baselineskip}
}
\newenvironment{inscriptionblock}{%
  \par\medskip
  \noindent\begin{minipage}{\textwidth}
  \setlength{\parindent}{0pt}
}{%
  \end{minipage}
  \medskip
}

\title{Oracle Inscription Test --- Original Book Format}
\begin{document}
\maketitle

\section*{Example: Paired Verification}

\inscriptionlabel{Preface:} \inscriptionlabel{Charge:} 
\inscriptionlabel{Prognostication:} \inscriptionlabel{Verification:}

\inscriptionsection{PAIR OF COMPLEMENTARY CHARGES}

\begin{inscriptionblock}
\inscriptionside{positive, left side}
\inscriptionref{Ping-pien 235.2}

Crack-making on chi-mao (day 16), Ch'üch divined:
``It will rain.''
The king, reading the cracks (said):
``It will rain; it will be a jen day.''
On jen-wu (day 19) it really did rain.
\end{inscriptionblock}

\begin{inscriptionblock}
\inscriptionside{negative, right side}
\inscriptionref{Ping-pien 235.1}

Crack-making on chi-mao (day 16), Ch'üch divined:
``It will not perhaps rain.''
\end{inscriptionblock}

In period I, verifications (and prognostications) might provide considerable detail; for example:

\inscriptionlabel{Preface:} \inscriptionlabel{Charge:}
\inscriptionlabel{Prognostication:}

\inscriptionsection{PAIR OF COMPLEMENTARY CHARGES}

\begin{inscriptionblock}
\inscriptionside{negative, left side}
\inscriptionref{Ping-pien 1.4}

Crack-making on kuei-ch'ou (day 50), Cheng divined:
``From today to ting-ssu (day 54) we will not perhaps harm the Chou.''
\end{inscriptionblock}

\begin{inscriptionblock}
\inscriptionside{positive, right side}
\inscriptionref{Ping-pien 1.3}

Crack-making on kuei-ch'ou (day 50), Cheng divined:
``From today to ting-ssu (day 54) we will harm the Chou.''
The king, reading the cracks, said:
``(Down to) ting-ssu (day 54) we should not perhaps harm (them); on the coming chia-tzu (day 1) we will harm (them).''
On the eleventh day, kuei-hai (day 60), (our) chariots did not harm (them); in the t'ou period between that evening and chia-tzu (day 1), (we) really harmed (them).
\end{inscriptionblock}

\end{document}
